\subsection{FusedLoRA and FusedMultiLoRA}
\label{sec:fused-lora-and-fused-multi-lora}

As described in Section~\ref{sec:motivation-lora-overhead}, LoRA modules introduce significant runtime overhead despite adding only a small number of parameters. 
Our profiling reveals that this overhead stems primarily from redundant memory access. 
Our goal is to fuse operations within the LoRA computation graph to reduce these memory transfers while preserving compute efficiency.

\myparagraph{Design Considerations and Challenges}
\begin{sloppypar}
While kernel fusion can reduce redundant memory access, fusing all LoRA operations into a single kernel introduces practical challenges.
First, the frozen GEMM operation $Y_1 = XW$ is compute-bound and highly sensitive to kernel tiling strategies and GPU resource usage (e.g., shared memory and register file). 
A suboptimal tiling layout or overuse of registers and shared memory can greatly degrade its performance.
Second, fusing operations with producer-consumer dependencies, such as $\widehat{X}A$ followed by $(\widehat{X}A)B$, may require recomputing intermediate results or introducing thread block synchronization, both of which can add overhead~\cite{shi2023welder}.
Thus, a good fusion strategy must minimize memory access while preserving optimal compute performance and avoiding expensive synchronization or recomputation.
\end{sloppypar}

\begin{figure}[t]
    \centering
    \includegraphics[width=0.99\linewidth]{figures/system-design-fig-2-fusion-design.pdf}
    \caption{Overview of our fusion strategy for LoRA modules in the forward pass, illustrating the full graph fusion approach vs. the split graph fusion approach.}
    \Description{Fusion design.}
    \label{fig:system-design-fig-2-fusion-design}
\end{figure}


\myparagraph{Full Graph Fusion vs. Split Graph Fusion} 
Figure~\ref{fig:system-design-fig-2-fusion-design} illustrates three design choices for handling the intermediate tensor $S = \hat{X}A$ in the forward pass. 
The first option recomputes $S$ inside each tile of the fused kernel, but requires loading the entire $A$ matrix repeatedly, becoming expensive when batch size $M$ is large. 
The second option fuses computation and uses synchronization across thread blocks to share $S$, where only a single $Mtile$ (the m-th block in the token dimension) computes the intermediate $S$ tiles and writes to global memory, while other tiles wait through a semaphore. This adds coordination overhead. 
Our approach takes a third option: explicitly storing and reloading $S$ from GPU global memory. Since $S$ is much smaller than other tensors and depends on the small LoRA rank $r$, the cost of reading and writing it is low. 
Splitting the graph at $S$ avoids both recomputation and synchronization while reducing expensive memory traffic associated with full-sized activation tensors. 
This approach preserves GPU resources for optimal tiling of the compute-bound $XW$ operation, maintaining peak performance for the most computationally intensive part.



\begin{figure*}[t]
    \centering
    \includegraphics[width=0.99\linewidth]{figures/system-design-fig-1-lora-kernel.pdf}
    \caption{LoRA kernel design. FusedLoRA reduces memory accesses by combining memory-heavy LoRA branches with base GEMM operations on frozen weights. The right figure has transposed weight tensors to match the hardware memory layout.}
    \label{fig:system-design-fig-1-lora-kernel}
\end{figure*}

\myparagraph{FusedLoRA Design}
Figure~\ref{fig:system-design-fig-1-lora-kernel} illustrates our fused kernel design for both forward and backward passes.
In the forward pass (Figure~\ref{fig:system-design-fig-1-lora-kernel}(a)), 
we combine dropout and down-projection into a single kernel (\filledone) to eliminate reloading of the full-sized activation tensor.
We also fuse the compute-bound base model GEMM ($Y_1 = XW$) with the memory-bound LoRA operations ($Y_2 = \alpha SB$) (\filledtwo). 
This fusion eliminates redundant memory operations and saves one read and write of the full-sized output tensor by directly accumulating partial results, without affecting the base GEMM performance.
In the backward pass, we apply similar principles. Operation \filledthree fuses the gradient computation of $S$ and $B$, eliminating the need to reload $dY$.
Operation \filledfour remains separate since it operates on the small masked input, where fusion provides minimal benefit.
Operation \filledfive horizontally fuses the compute-intensive gradient computation for the base model with memory-bound LoRA path operations, preventing redundant reads and writes of partial output gradients.
The key insight of our approach is strategically identifying operations where horizontal fusion reduces memory traffic without compromising computational efficiency. By fusing operations that share large tensors (\filledtwo, \filledthree, and \filledfive), we significantly reduce memory bottlenecks while maintaining optimal tiling strategies for compute-bound operations.

\begin{figure}[t]
    \centering
    \includegraphics[width=0.975\linewidth]{figures/system-design-fig-3-fused-multi-lora.pdf}
    \caption{Illustration of FusedMultiLoRA in the forward pass. The routing of LoRA adapters is done at the tile level.}
    \Description{FusedMultiLoRA.}
    \label{fig:system-design-fig-3-fused-multi-lora}
\end{figure}


\myparagraph{Extending to FusedMultiLoRA}
To support concurrent fine-tuning of multiple LoRA adapters, we extend FusedLoRA to FusedMultiLoRA, allowing our fused kernels to operate on mixed-adapter batches from different jobs. 
As shown in Figure~\ref{fig:system-design-fig-3-fused-multi-lora}, each input $Mtile$ is tagged with an adapter ID and configuration, such as LoRA rank, scaling factor, and dropout ratio, stored in a lightweight lookup table. 
During execution, the frozen model computation is shared across all tokens, while adapter-specific logic is applied dynamically per $Mtile$. 
For each ($Mtile, Ntile$) tile of the output, the kernel loads the appropriate $A$ and $B$ matrices and applies the correct scaling and dropout. 
In the backward pass, the same mapping is used to route gradients to their respective adapters without interference. 
This tile-level routing allows efficient execution of heterogeneous adapters in a single fused run, avoiding redundant computation and enabling the job-level optimizations introduced in Section~\ref{sec:scheduler}.


FusedLoRA reduces memory traffic by fusing LoRA operations around shared activations while preserving base model efficiency. Building on FusedLoRA, FusedMultiLoRA supports heterogeneous adapters via tile-level routing. The system dynamically chooses between them, falling back to FusedLoRA when only one adapter is present in the batch.